Petroleum statistics classify aviation fuel by customer, whereas greenhouse-gas inventories
classify flights by origin and destination. We develop a source-row-conserving bridge between
these systems for Pakistan using annual fuel reports, five matched fiscal years of OECD
aviation estimates and 42 monthly comparisons. In FY2023--24, commercial aviation fuel implies
\SI{2.032}{\mega\tonne} \COtwo{}, 0.8\% below OECD's combined domestic and international
territory estimate. This aggregate agreement is misleading: resident-customer fuel is 28.6\%
above OECD's resident-carrier estimate, while foreign-customer fuel is 23.1\% below its
non-resident estimate, causing 97\% of the component discrepancy to cancel.

The figure the national inventory carries is the aggregate movement model's
\SI{0.301}{\mega\tonne} \COe{} (0.299 Mt as \COtwo{} alone) --- the estimator with the
strongest activity evidence behind it, and the one every national product uses. Around it:
the OECD resident-carrier partition gives \SI{0.292}{\mega\tonne} of domestic \COtwo{},
which is a calibrated consistency check rather than independent validation because OECD
supplies that partition, and the route--aircraft Monte Carlo gives a mean of 0.244 Mt
\COtwo{} with a 95\% interval of 0.215--0.272 Mt. Pakistan's official 2022 reporting instead applies an 80/20
domestic--international allocation that is not consistent with the petroleum and
flight-geography evidence examined here. The framework conserves every reported fuel row,
keeps unresolved quantities explicit and separates calibration, external evaluation,
structural uncertainty and parametric uncertainty. It provides a transferable method for
reconciling purchaser-based petroleum statistics with flight-based aviation accounts.

The common annual core now reaches FY2024--25. Its LTO allocation uses the corresponding complete PCAA airport-movement source, with the reference aircraft composition explicitly carried where own-year fleet evidence is absent. The detailed route, OECD and structural comparisons in this note remain anchored to FY2023--24.
